% !TEX root = ../Projektdokumentation.tex
\section{Projektplanung} 
\label{sec:Projektplanung}
Um einen strukturierten Ablauf des Projektes zu gewährleisten, werden in diesem Abschnitt
die Projektphasen, Abweichungen vom Projektantrag, benötigte Ressourcen sowie der Entwicklungsprozess beschrieben.

\subsection{Projektphasen}
\label{sec:Projektphasen}
Die Umsetzung des Projektes erfolgt innerhalb von zwei vierwöchigen Sprints.
Gemäß den Vorgaben der IHK wird das Projekt, wie in Tabelle~\ref{tab:Zeitplanung} dargestellt,
innerhalb von 80 Stunden bearbeitet.

\paragraph{Zeitplanung}
Tabelle~\ref{tab:Zeitplanung} zeigt die grobe Zeitplanung.
\tabelle{Zeitplanung}{tab:Zeitplanung}{ZeitplanungKurz}\\
Eine detailliertere Zeitplanung findet sich im \Anhang{app:Zeitplanung}.

\subsection{Abweichungen vom Projektantrag}
\label{sec:AbweichungenProjektantrag}
Im Vergleich zum Projektantrag wurde die geplante Anbindung von Hangfire durch die Anbindung an RabbitMQ ersetzt.
Während der Planung wurde festgestellt, dass Hangfire für die erste Iteration nicht benötigt wird.
Da die Synchronisation ereignisbasiert erfolgen soll, wurde stattdessen RabbitMQ als Message Broker vorgesehen.

\subsection{Ressourcenplanung}
\label{sec:Ressourcenplanung}
Für die Durchführung des Projektes werden ein Büroarbeitsplatz mit Internetzugang,
ein Entwicklungsgerät mit zusätzlichen Bildschirmen sowie Serverressourcen für Anwendung und Datenbank benötigt.
Als Software kommen Windows 11, JetBrains Rider 2026, RabbitMQ, PostgreSQL 16,
Debian 12.13 sowie HeidiSQL und pgAdmin 4 zur Datenbankverwaltung zum Einsatz.
\par\smallskip
Zusätzlich werden während der Review- und Testphase zwei weitere Entwickler zur Qualitätssicherung einbezogen.

\subsection{Entwicklungsprozess}
\label{sec:Entwicklungsprozess}
Während der Umsetzung wird ein agiler Entwicklungsprozess nach Scrum verwendet.
Das Projekt wird dabei in kleinere Arbeitspakete aufgeteilt, die innerhalb der Sprints umgesetzt werden.
Dadurch kann flexibel auf technische Probleme oder geänderte Anforderungen reagiert werden.
\par\smallskip
Die Planung und Nachverfolgung der Aufgaben erfolgt über Tickets auf einem Jira-Board.
Die Tickets beschreiben die einzelnen Arbeitspakete und dienen während des Projektes zur Fortschrittskontrolle.
Eine Übersicht der verwendeten Tickets befindet sich im \Anhang{app:Tickets}.